\documentclass[11pt,a4paper]{article}
\usepackage{ctex}
\usepackage{amsmath}
\usepackage{amssymb}
\usepackage{geometry}
\geometry{left=2.5cm,right=2.5cm,top=2.5cm,bottom=2.5cm}
\begin{document}

\title{2025年春季学期大学物理学期中考试（修正版）}
\author{整理：廖敬远}
\date{}
\maketitle

\section*{一、填空题（每题2分，共18分）}

\begin{enumerate}
    \item 质点沿直线运动，加速度满足 \( a = -kx \)。初始条件：\( t=0 \) 时，\( x=0 \)，\( v=v_0 \)。求速度 \( v \) 关于位置 \( x \) 的表达式。

    \item 小明以 \( v=3\ \mathrm{m/s} \) 向东奔跑时，感觉风从正南方吹来；当他以 \( 6\ \mathrm{m/s} \) 向东奔跑时，感觉风从正东南方向吹来。求实际风速的大小。

    \item 光滑水平面上，距转轴分别为 \( r \) 和 \( 2r \) 处各有一质量为 \( m \) 的物块，转轴与物块、物块与物块之间均用轻弹簧相连。当系统以恒定角速度 \( \omega \) 转动且两弹簧长度相等时，两弹簧弹力大小之比为多少？
\vspace{0.5em}

\textbf{解：}

\noindent \textbf{1. 模型与受力分析}
\begin{itemize}
    \item 设转轴与物块1（距转轴 \( r \) 处）之间的弹簧弹力为 \( F_1 \)，物块1与物块2（距转轴 \( 2r \) 处）之间的弹簧弹力为 \( F_2 \)。
    \item 由“两弹簧长度相等”的条件，可知弹簧1长度为 \( r \)，弹簧2长度为 \( 2r - r = r \)，两弹簧均处于拉伸状态，弹力为拉力。
    \item 规定指向转轴的方向为正方向，向心力的合力沿正方向。
\end{itemize}

\noindent \textbf{2. 对物块2列向心力方程}
物块2做匀速圆周运动，半径为 \( 2r \)，向心力仅由弹簧2的拉力提供：
\[
F_2 = m \omega^2 \cdot 2r
\tag{1}
\]

\noindent \textbf{3. 对物块1列向心力方程}
物块1做匀速圆周运动，半径为 \( r \)，受两个弹簧的拉力：弹簧1向内的拉力 \( F_1 \)（正方向）、弹簧2向外的拉力 \( F_2 \)（负方向），合力提供向心力：
\[
F_1 - F_2 = m \omega^2 \cdot r
\tag{2}
\]

\noindent \textbf{4. 联立求解弹力比}
将式(1)代入式(2)：
\[
F_1 - 2m\omega^2 r = m\omega^2 r
\]
整理得：
\[
F_1 = 3m\omega^2 r
\]

因此，两弹簧弹力大小之比为：
\[
\frac{F_1}{F_2} = \frac{3m\omega^2 r}{2m\omega^2 r} = \frac{3}{2}
\]
    \item 一颗子弹射向三个并排紧靠的相同木块，子弹在每个木块中受到的阻力大小相同，穿过每个木块的时间也相同。求子弹穿出后，三个木块从左到右的最终速度之比。

    \item 倾角为 \( \theta \) 的光滑斜面上，一滑块与斜面顶端的弹簧相连，弹簧劲度系数为 \( k \)，初始时滑块静止。现给滑块沿斜面向下的初速度 \( v_0 \)，当弹簧伸长量为 \( x \) 时，求滑块的动能大小。

    \item 电梯内有一质量为 \( m \) 的物体，相对电梯静止。电梯以竖直向下的加速度 \( a = \dfrac{g}{4} \) 向下运动高度 \( h \)，求此过程中电梯对物体做的功。

    \item 一空心圆柱底面为圆环形，内半径 \( R_1 \)，外半径 \( R_2 \)，总质量为 \( m \)，且质量均匀分布。求该柱体绕中心几何轴线的转动惯量。

    \item 某不稳定粒子静止质量为 \( m_0 \)，固有衰变时间为 \( \tau_0 \)，在实验室参考系中测得衰变时间为 \( \tau \)。求实验室中测得该粒子的动能。

\vspace{0.5em}
\textbf{解：}
\begin{enumerate}
    \item \textbf{时间膨胀效应求洛伦兹因子}
    粒子的固有衰变时间 \( \tau_0 \) 为原时，实验室参考系中时间膨胀，满足公式：
    \[
    \tau = \frac{\tau_0}{\sqrt{1-\dfrac{v^2}{c^2}}} = \gamma \tau_0
    \]
    其中洛伦兹因子 \( \gamma = \dfrac{1}{\sqrt{1-\dfrac{v^2}{c^2}}} \)，整理得：
    \[
    \gamma = \frac{\tau}{\tau_0}
    \]

    \item \textbf{相对论动能公式}
    相对论中，粒子的动能等于总能量与静能的差值，公式为：
    \[
    E_k = \gamma m_0 c^2 - m_0 c^2 = (\gamma - 1) m_0 c^2
    \]

    \item \textbf{代入求解最终动能}
    将 \( \gamma = \dfrac{\tau}{\tau_0} \) 代入动能公式：
    \[
    E_k = \left( \frac{\tau}{\tau_0} - 1 \right) m_0 c^2
    \]
\end{enumerate}
    \item 从上小下大的喷头中竖直向上喷出的水柱最大高度为 \( H \)，喷头出口处高度不计，喷头上截面半径为 \( r \)，下截面半径为 \( R \)，水的密度为 \( \rho \)，大气压强为 \( p_0 \)。求喷头下截面处水的压强。
\textbf{解：}
\begin{enumerate}
    \item 由机械能守恒验证：【取出口处为重力势能零点，一小段水的机械能守恒：出口动能=最高点重力势能】
    \[
    \frac{1}{2}\Delta m v_1^2 = \Delta m g H \implies v_1 = \sqrt{2gH}
    \]

    \item 连续性方程求下截面流速
    上截面面积 \( S_1 = \pi r^2 \)，下截面面积 \( S_2 = \pi R^2 \)，
    
    由 \( S_1 v_1 = S_2 v_2 \)：
    \[
    v_2 = \frac{r^2}{R^2} \sqrt{2gH}
    \]

     \item \textbf{伯努利方程（计入高度差 h）求下截面压强}
    
     设下截面为位置1（高度 \( 0 \)，压强 \( p \)，流速 \( v_2 \)），
    
     最高点为位置2（高度 \( h+H \)，压强 \( p_0 \)，流速 0。
    
     伯努利方程：

    合并化简：
     \( p_0 + \rho g \left( h + H\left(1 - \dfrac{r^4}{R^4}\right) \right) \)

\end{enumerate}
\end{enumerate}

\section*{二、计算题（共32分）}

\begin{enumerate}
    \item（10分）光滑水平面上有一长为 \( L \)、质量均匀分布为 \( m \) 的细杆，以垂直于杆方向的平动速度 \( v_0 \) 运动。杆前方有一固定立柱 \( OO' \)，细杆与立柱发生完全非弹性碰撞，碰撞点距杆一端 \( \dfrac{L}{3} \)。求碰撞瞬间细杆的角速度大小。
\vspace{0.5em}

\textbf{解：}
\begin{enumerate}
    \item \textbf{角动量守恒条件分析}
    碰撞过程时间极短，立柱对杆的作用力通过碰撞点，对碰撞点的力矩为零，因此细杆对碰撞点的角动量守恒。

    \item \textbf{计算碰撞前杆对碰撞点的角动量}
    杆做平动，质心位于杆中点，距离碰撞点的距离为：
    \[
    d = \frac{L}{2} - \frac{L}{3} = \frac{L}{6}
    \]
    杆的质心速度为 \( v_0 \)，方向垂直于杆，因此对碰撞点的角动量大小为：
    \[
    L_1 = m v_0 d = m v_0 \cdot \frac{L}{6}
    \]

    \item \textbf{计算碰撞后杆绕碰撞点的转动惯量}
    细杆对质心的转动惯量为：
    \[
    I_{\text{cm}} = \frac{1}{12} m L^2
    \]
    由平行轴定理，杆绕碰撞点的转动惯量为：
    \[
    I = I_{\text{cm}} + m d^2 = \frac{1}{12} m L^2 + m \left( \frac{L}{6} \right)^2
    \]
    化简得：
    \[
    I = \frac{1}{12} m L^2 + \frac{1}{36} m L^2 = \frac{1}{9} m L^2
    \]

    \item \textbf{角动量守恒方程求解角速度}
    碰撞后杆绕碰撞点转动，角动量为 \( L_2 = I \omega \)。
    由角动量守恒 \( L_1 = L_2 \)，得：
    \[
    m v_0 \cdot \frac{L}{6} = \frac{1}{9} m L^2 \cdot \omega
    \]
    约去 \( m \) 并整理得：
    \[
    \omega = \frac{3 v_0}{2 L}
    \]
\end{enumerate}
    \item（10分）光滑水平桌面上，左车、轻弹簧、松弛轻绳、轻弹簧、右车依次连接。左车质量 \( m_1 \)，右车质量 \( m_2 \)，两弹簧劲度系数分别为 \( k_1 \)、\( k_2 \)。初始时整个系统静止，给左车一瞬时向左的初速度 \( v_0 \)。求当两车达到共同速度时，两弹簧的伸长量 \( \Delta x_1 \) 和 \( \Delta x_2 \)。

    \item（6分）在惯性系 \( S \) 中，事件1和事件2发生在同一地点，事件1发生后经过 \( 3\ \mu\mathrm{s} \) 事件2发生。在惯性系 \( S' \) 中观测，两事件的时间间隔为 \( 5\ \mu\mathrm{s} \)。求 \( S' \) 系中两事件发生的位置间隔大小。
\end{enumerate}

\end{document}